\chapter*{Аннотация}
\addcontentsline{toc}{chapter}{Аннотация}

Диссертация посвящена послойному и вычислительно сопоставимому исследованию
режимов предиктивного кодирования (PC) относительно обратного распространения
ошибки (BP). Работа разделяет четыре аналитических измерения, которые часто
смешиваются при сравнении алгоритмов обучения: сходство конечного поведения,
сходство внутреннего механизма, вычислительную стоимость и возможность
до штатного завершения заменить оставшуюся каноническую итеративную часть
зарегистрированным ограниченным аналитическим завершением при сохранении
требуемого ответа. Для перехода
от наблюдения внутреннего состояния к вычислительному решению вводятся две
ограниченные теоретические конструкции: PC-TREF задаёт эквивалентность
состояний относительно требуемого действия, полную достаточность
диагностического представления на зарегистрированной области и более узкую
селективную достаточность правила
для конкретного действия, а PC-CATM описывает формирование, перенос и
компенсацию каналов коррекции. Архитектура QWake-PC операционализирует эти основания; в работе
экспериментально проверяется её ограниченная реализация QWake-FP для FixedPred.
В QWake-FP раннее действие означает выбор зарегистрированного аналитического
завершения \nolinkurl{fixedpred_eta1_wavefront_completion_v1} вместо
оставшихся канонических итеративных проходов; полный точный остаток
\nolinkurl{complete_suffix_stage2_baseline_v1} служит эталонным и
резервным путём.

Экспериментальная программа организована как последовательность заранее
зафиксированных этапов. Этапы этап~1 и этап~2 сравнивают BP, Exact, FixedPred
и Strict в контролируемых условиях. этап~3A исследует послойные градиенты и
представления; этап~3B локализует стоимость итеративного вывода состояния и
проверяет кандидаты точного вычислительного замещения посредством заранее
зарегистрированных проверок эквивалентности. Финальная ветка QWake-FP разделяет
три последовательных вопроса: присутствуют ли на поверхности траекторий с
зафиксированной целостностью предтерминальные записи, состояния которых допускают
раннее действие относительно задачи; можно ли на зафиксированной калибровочной
поверхности селективно распознавать ненулевое подмножество таких записей по состоянию,
доступному до действия, без наблюдавшихся опасных принятий; и даёт ли такое
решение положительную выгоду при
заранее зафиксированном полном учёте стоимости принятия решения.

В исследованных условиях FixedPred демонстрировал более высокую наблюдаемую
близость к BP по направлению градиента и внутренним представлениям, чем Strict,
при одновременном уменьшении нормы
градиента в ранних слоях. Профилирование этап~3B связало существенную часть
стоимости исследованных режимов PC с итеративным выводом состояния; последующие
кандидаты B1/B2 прошли зарегистрированные проверки эквивалентности, а полная
матрица сопоставленного профилирования была завершена без повторных попыток.
При этом зарегистрированный инженерный экран продолжения дал 0/16
квалифицированных конфигураций во всех четырёх группах B1/B2×FixedPred/Strict,
и оба кандидата получили решение \texttt{reject\_or\_revise}; этот результат не
интерпретируется как утверждение о превосходстве исходной линии. В QWake-FP поверхность с зафиксированной целостностью содержала предтерминальные
записи, состояния которых эталонная проверка относила к классу
\texttt{EARLY\_ADMISSIBLE}, а ненулевая селективная распознаваемость их
подмножества по состоянию до действия с нулём наблюдавшихся опасных принятий
наблюдалась внутри замороженного семейства из 2625 скалярных правил принятия
решения на 756 калибровочных записях. Из них 264 правила имели ненулевое
покрытие при отсутствии опасных принятых решений, а правило с наибольшим
покрытием среди правил с нулём наблюдавшихся опасных принятий покрывало 216
из 756 калибровочных записей (28.57\%). Само правило использует только
\texttt{compute\_step >= 5}; поэтому положительный результат C08 на этой
поверхности задаёт временную границу, эквивалентную фиксированному префиксу, а
не демонстрирует зависимость решения от содержания входа или преимущество
механизмных признаков PC-CATM. Структурная сверка временной поверхности
показывает, что эти 216 записей состоят из 108 предтерминальных записей шага 5 с
одним оставшимся проходом и 108 записей терминальной границы шага 6; поэтому
28.57\% сохраняется как зарегистрированное покрытие полной поверхности и не
трактуется как доля предтерминальных применений раннего аналитического действия. При этом ни одно правило
не обеспечило положительную совокупную чистую экономию при зафиксированном
полном учёте стоимости принятия решения. В полной стоимости этого правила
преобладали накладные расходы наблюдателя; правило C2 не было зафиксировано,
поэтому подтверждающий этап C3 не открывался.

Полученный отрицательный экономический результат относится именно к полному
учёту стоимости принятия решения, зафиксированному для C2, и не измеряет
добавочную стоимость выполнения минимального распознавателя. Поэтому работа
поддерживает вывод о наличии в исследованной области предтерминальных записей,
состояния которых допускают раннее действие относительно задачи, и ненулевой
селективной калибровочной распознаваемости их подмножества по состоянию до
действия с нулём наблюдавшихся опасных принятий;
отклоняет гипотезу о существовании
правила с нулём наблюдавшихся опасных принятий и положительной экономической
выгодой в замороженном семействе C2 при
указанном учёте стоимости; и оставляет экономическую осуществимость минимальной
реализации отдельной непроверенной оцениваемой величиной.

\noindent\textbf{Ключевые слова:} предиктивное кодирование, обратное
распространение ошибки, послойный анализ, вычислительная стоимость,
аналитическое завершение, трассируемость.

\clearpage
\begin{otherlanguage}{english}
\chapter*{Abstract}
\addcontentsline{toc}{chapter}{Abstract}

This dissertation studies layer-wise and computationally controlled
comparisons of predictive-coding (PC) regimes with backpropagation (BP). The
research separates four analytical dimensions that are often conflated:
similarity of final behavior, similarity of internal mechanism, computational
cost, and the possibility of terminating part of the computation before
its normal completion subject to a registered dangerous-accept criterion. To connect internal diagnostics with a computational
decision, the dissertation introduces two bounded theoretical constructs:
PC-TREF defines task-relative state equivalence, a full sufficiency criterion
on the registered state domain, and a narrower action-selective
sufficiency condition, while PC-CATM models correction formation, transport,
and cancellation.
QWake-PC operationalizes these foundations as a compute-control architecture;
only its bounded FixedPred instantiation, QWake-FP, is evaluated empirically.
In QWake-FP, the early action replaces the remaining canonical iterative suffix
with the registered bounded analytic completion
\nolinkurl{fixedpred_eta1_wavefront_completion_v1}; the complete exact
suffix \nolinkurl{complete_suffix_stage2_baseline_v1} remains the
reference and fallback path.

The experimental program is organized as a sequence of preregistered and
versioned stages. Stage~1 and Stage~2 compare BP, Exact, FixedPred, and Strict
under controlled conditions. Stage~3A examines layer-wise gradients and
representations, while Stage~3B localizes state-inference cost and evaluates
exact replacement candidates through registered equivalence gates. The final
QWake-FP branch separates three questions: whether the sealed trajectory surface
contains preterminal records whose states make the registered analytic-completion action task-admissible,
whether a non-zero subset can be selectively recognized from the pre-action
state with zero observed dangerous accepts on the frozen calibration surface,
and whether such a decision yields
positive savings under frozen full decision-cost accounting.

Within the investigated conditions, FixedPred showed higher observed similarity
to BP than Strict in gradient direction and learned representations, while its
early-layer gradient
norm was substantially reduced. Stage~3B profiling associated a substantial
share of the measured PC cost with state inference; subsequent B1/B2 candidates
passed their registered equivalence gates, and the matched profiling matrix was
completed without retries. The registered engineering-continuation screen then
qualified 0 of 16 configurations in each of the four B1/B2-by-FixedPred/Strict
groups, and both candidates received \texttt{reject\_or\_revise}; this bounded
negative result is not interpreted as a superiority claim for the baseline. In QWake-FP, the sealed surface contained preterminal records whose states
made the registered analytic-completion action task-admissible, and non-zero selective pre-action
recognizability with zero observed dangerous accepts was found within the frozen
family of 2,625 scalar rules on 756 calibration records: 264 rules had
non-zero coverage with zero dangerous accepts, and the highest-coverage rule
among rules with zero observed dangerous accepts covered 216 of 756
calibration records (28.57\%). That rule uses only
\texttt{compute\_step >= 5}; therefore C08 establishes a temporal boundary
that is equivalent to a fixed prefix on this surface, not input-dependent
adaptivity or superiority of PC-CATM features. Structural reconciliation of the
temporal surface shows that these 216 records comprise 108 preterminal records
at step 5 with one remaining sweep and 108 records at the terminal boundary
step 6; therefore 28.57\% remains the registered full-surface coverage and is
not interpreted as preterminal analytic-action coverage. No rule achieved positive
aggregate net saving under the frozen full decision-cost accounting. The full
cost assigned to that highest-coverage zero-danger rule was overwhelmingly
observer-dominated, no C2 rule freeze was established, and confirmatory C3
therefore remained closed.

The negative economic result is specific to the frozen full-cost accounting
used in C2 and does not measure the marginal runtime cost of a minimal
recognizer. The dissertation therefore supports the presence of preterminal records whose
states make the registered analytic-completion action task-admissible and non-zero selective pre-action
recognizability with zero observed dangerous accepts on the frozen calibration
surface, rejects the existence of a rule with zero observed dangerous accepts,
non-zero coverage, and positive aggregate net saving in the frozen C2 family
under that accounting
regime, and leaves marginal implementation viability as a separate untested
estimand.

\noindent\textbf{Keywords:} predictive coding, backpropagation, layer-wise
analysis, computational cost, analytic completion, traceability.
\end{otherlanguage}
